% SPDX-License-Identifier: CC-BY-4.0
% Copyright (C) 2026 Fabian Januszewski, Abhijit (the Langlandscape team)
% Definitions and remarks of the chapters (mathematics/chapters/s0N) that Sections 1--5 refer to by label.
% Each entry: \rptnum{main.pdf number}\label{rpt:<chapter label>} + bold run-in heading. Numbers are those of main.pdf.
% Definitions are verbatim where short, otherwise marked "(condensed)"; statuses are copied, never upgraded.

\noindent The following entries reproduce, in the numbering of \mainpdf, the definitions, propositions, remarks and chapter sections
that the preceding sections cite. Short definitions are quoted; longer ones are condensed and marked as such. Status markers are those
of the chapters.

\subsubsection*{Chapter 1: classical objects}

\noindent\rptnum{1.1}\label{rpt:def:s01-B}\textbf{Definition 1.1 (definite algebra).}
$D>1$ is squarefree with $\omega(D)$ odd, and $B=\Bd$ is the quaternion algebra over $\Q$ ramified exactly at $\infty$
and the primes $\ell\mid D$. $\Nrd$ is the reduced norm, $B^1=\ker\Nrd$; since $B$ is definite, $\Nrd(B^\times)=\Q_{>0}$.
$\cO\subset B$ is a maximal order. For $\ell\nmid D$ we fix $\iota_\ell:\cO_\ell\xrightarrow{\sim}M_2(\Z_\ell)$; for $\ell\mid D$,
$\cO_\ell=\{x\in B_\ell:\Nrd x\in\Z_\ell\}$ is the unique maximal order of the division algebra $B_\ell$, with a
uniformiser $\Pi_\ell$, $\Pi_\ell^2\in\ell\,\cO_\ell^\times$.
Throughout, $\Zhat=\prod_\ell\Z_\ell$, $\hat\Q=\Zhat\otimes\Q$ (finite adeles), $B_\ell=B\otimes\Q_\ell$, $\cO_\ell=\cO\otimes\Z_\ell$,
$\hat B=B\otimes\hat\Q$, $\hat\cO=\cO\otimes\Zhat=\prod_\ell\cO_\ell$.
\stknown{} (existence and uniqueness of $B$ up to isomorphism, local maximal orders: \cite[Thm.~13.3.11, p.~206; Prop.~13.3.4, p.~205]{Vo}).

\medskip
\noindent\rptnum{1.2}\label{rpt:def:s01-levels}\textbf{Definition 1.2 (levels).}
For integers $N\ge1$ with $(N,D)=1$ and matrices read through $\iota_\ell$ (so $\ell\nmid D$ whenever $\ell\mid N$):
$\Uzero{N}=\prod_\ell U_{0,\ell}$, $\Uone{N}=\prod_\ell U_{1,\ell}$ with $U_{\bullet,\ell}=\cO_\ell^\times$ for $\ell\nmid N$ and
\[U_{0,\ell}=\{\left(\begin{smallmatrix}a&b\\c&d\end{smallmatrix}\right)\in\GL_2(\Z_\ell):c\equiv0\},\qquad
U_{1,\ell}=\{\left(\begin{smallmatrix}a&b\\c&d\end{smallmatrix}\right)\in\GL_2(\Z_\ell):c\equiv0,\ d\equiv1\}\pmod{\ell^{v_\ell(N)}}.\]
Here $N=1$ gives $\Uzero1=\Uone1=\hat\cO^\times$. Further $U_s:=\Uone{Mp^s}$ ($M\ge1$, $s\ge0$, $p$ prime, $(Mp,D)=1$), $S_s:=\XB{U_s}$,
$\Delta_s:=(1+p\Zp)/(1+p^s\Zp)$ for $s\ge1$, and for $q\nmid D$ and $U$ with $U_q=\cO_q^\times$: $U_0(q)U:=U\cap\Uzero{q}$ (i.e.\ $U_q$
replaced by $\Iw_q$).
(The chapter adds: levels start at $N\ge1$; $\Uone N\subseteq\Uzero N\subseteq\hat\cO^\times$, $\Uone{N'}\subseteq\Uone N$ for $N\mid N'$,
$-1\in\Uzero N$ always and $-1\in\Uone N$ iff $N\in\{1,2\}$. The \emph{lower} row is $\equiv(0,1)$, the convention of
\cite[\S4, p.~10]{Buz04}; it matters for $U_p$, Lemma~\ref{lem:s01-Up}.)

\medskip
\noindent\rptnum{1.4}\label{rpt:prop:s01-eichler}\textbf{Proposition 1.4 (Eichler orders).}
\emph{For $(N,D)=1$, $\Uzero{N}=\hat\cO(N)^\times$ for the Eichler order $\cO(N)=B\cap\hat\cO(N)$ of level $N$, locally
$\left(\begin{smallmatrix}\Z_\ell&\Z_\ell\\N\Z_\ell&\Z_\ell\end{smallmatrix}\right)$.} \stknown{}
\cite[Def.~23.4.1 and Prop.~23.4.3, p.~382]{Vo}.
(The pre-existing Lean material has no $\Uone{N}$, no $p$-tower and no Iwahori level at a prime outside the Eichler level: these are new.)
\begin{outline}
(Sketch of the identification, as in the chapter.) At $\ell\nmid N$ both sides are $\cO_\ell^\times$. At $\ell\mid N$ (so $\ell\nmid D$) the
local order $\iota_\ell^{-1}\left(\begin{smallmatrix}\Z_\ell&\Z_\ell\\N\Z_\ell&\Z_\ell\end{smallmatrix}\right)$ is the intersection of
$\cO_\ell$ with its conjugate by $\mathrm{diag}(1,\ell^{v_\ell(N)})$, hence Eichler by Voight's Def.~23.4.1 and Prop.~23.4.3. An element of
this order is a unit iff its determinant is in $\Z_\ell^\times$; since $\det\equiv ad\pmod{\ell^{v_\ell(N)}}$, this means $a,d\in\Z_\ell^\times$,
i.e.\ the element lies in $\GL_2(\Z_\ell)$ with $c\equiv0$, which is $U_{0,\ell}$.
\end{outline}

\noindent\rptnum{1.5}\label{rpt:def:s01-classset}\textbf{Definition 1.5 (class set).}
For compact open $U\subseteq\hat\cO^\times$: $\XB{U}:=B^\times\backslash\Bhat^\times/U$,
$\Gamma_x:=B^\times\cap xUx^{-1}$ ($x\in\Bhat^\times$), and $w_x:=|\Gamma_x/(\Gamma_x\cap\{\pm1\})|$.
Thus $w_x=|\Gamma_x|/2$ if $-1\in U$ (e.g.\ $\Uzero{N}$, and $U_0(q)U$ whenever $-1\in U$) and $w_x=|\Gamma_x|$ if $-1\notin U$
(e.g.\ $\Uone{N}$, $N\ge3$).

\medskip
\noindent\rptnum{(1)}\label{rpt:eq:s01-Vk}\textbf{Equation (1) (the coefficient module $V_k$).}
$V_k=\Sym^{k-2}(\Zp^2)$, $n=k-2$, realised as homogeneous polynomials $P(X,Y)$ of degree $n$, with the \emph{left} action of the monoid
$M_2(\Zp)\cap\GL_2(\Qp)$
\[(g\cdot P)(X,Y)=P\big((X,Y)g\big)=P(aX+cY,\,bX+dY),\qquad g=\left(\begin{smallmatrix}a&b\\c&d\end{smallmatrix}\right).\tag{1}\]
It is a left action, integral on the whole monoid (needed for $U_p$) and continuous \stproved{}. The map $\mathrm{ev}:V_k\to\Zp$,
$P\mapsto P(0,1)$, satisfies $\mathrm{ev}(g\cdot P)\equiv d^n\,\mathrm{ev}(P)\pmod{c}$; this is the line through which Hida theory sees weight $k$.

\medskip
\noindent\rptnum{3.1}\label{rpt:def:s01-brandt}\textbf{Definition 3.1 (Brandt module).}
For a finitely generated $\Z$- or $\Zp$-module $V$ with continuous left action of $U_p$ (through $\iota_p$),
\[\Brandt{U}{V}:=\{f:\Bhat^\times\to V:\ f(\gamma xu)=u_p^{-1}f(x)\ (\gamma\in B^\times,\ u\in U)\}\ \cong\ \bigoplus_{[x]\in\XB{U}}V^{\,x_p^{-1}\Gamma_xx_p},\]
via $f\mapsto(f(x))_x$ over a set of representatives. For trivial $V=\Z$: $\Z^X:=\Brandt{U}{\Z}$ (functions, pullback),
$\Z[X]$ the free module on $X=\XB{U}$ (pushforward), $\Z[X]^0=\ker(\deg)$. Naive evaluation $\Z[X]\times\Z^X\to\Z$,
$([x],f)\mapsto f(x)$, is a perfect duality; every map below on $\Z[X]$ is the transpose of the same-named map on $\Z^X$.
\stknown{} \cite[(4.1)--(4.2) and Prop.~4.5, pp.~68--69]{Gr99} (there in the convention $f(\gamma g)=\gamma f(g)$); the isomorphism is
\stproved{} in the chapter.

\medskip
\noindent\rptnum{3.2}\label{rpt:def:s01-pairing}\textbf{Definition 3.2 (Brandt pairing and norm map).}
On $\Z[X]$: $\langle[x],[y]\rangle:=w_x\delta_{xy}$ (weights $w$; the imported Lean \texttt{weightedPairing} for $k=2$ uses $|\Gamma_x|$
instead). The norm map is $\nu:\Z[X]\to\Z^X$, $[x]\mapsto w_x\mathbf 1_x$; it is the map $a\mapsto\langle a,\cdot\rangle$ and is an
isomorphism over $\Zp$ whenever $p\nmid w_x$ for all $x$, in particular for $p\ge5$ (Lemma~\ref{lem:s01-stab}(c)).
For $V$ with $-1$ acting trivially, $\nu_V:\bigoplus_xV_{\Gamma_x}\to\bigoplus_xV^{\Gamma_x}$,
$v\mapsto\sum_{\gamma\in\Gamma_x/(\Gamma_x\cap\{\pm1\})}(x_p^{-1}\gamma x_p)\,v$ (for $V=\Z$ this is multiplication by $w_x$ on the summand $x$).
\stproved{}; \stformal{} (\lean{Langlandscape.Brandt.brandtPairing}, \lean{Langlandscape.Brandt.normMap},
\lean{Langlandscape.Brandt.normMap\_bijective}; $p$-adic version \lean{Langlandscape.Brandt.normMap\_bijective\_padic}).

\medskip
\noindent\rptnum{4.1}\label{rpt:def:s01-hecke}\textbf{Definition 4.1 (Hecke operators)} (condensed).
All levels are factorisable, $U=\prod_vU_v$; $\Sigma_p:=M_2(\Zp)\cap\GL_2(\Qp)$. For $g\in\Bhat^\times$ with $g_p\in\Sigma_p$ (via $\iota_p$; if
$V$ only carries a $U_p$-action, require $g_p\in U_p$) and $UgU=\coprod_ig_iU$ (a finite union):
\[(T_gf)(x)=\sum_ig_{i,p}\cdot f(xg_i)\qquad(f\in\Brandt{U}{V}).\]
On $\Z[X]$, $X=\XB{U}$, $T_g$ is the transpose for the naive evaluation pairing, $T_g[x]=\sum_i[xg_i]$ (trivial coefficients).
With $\varpi_\ell:=\iota_\ell^{-1}\mathrm{diag}(\ell,1)$ and $z_\ell$ the central idele $\ell$ at $\ell$ ($1$ elsewhere):
$T_\ell:=T_{\varpi_\ell}$ and $S_\ell:=T_{z_\ell}$ for $\ell\nmid DMpq$ ($q$ the auxiliary prime of a level $U_0(q)U$, if any);
$U_\ell:=T_{\varpi_\ell}$ for $\ell\mid M$; $S_\ell:=T_{z_\ell}$ also for $\ell\mid D$;
the diamonds $\langle d\rangle:=T_{\delta_d}$, $d\in(\Z/Mp^s)^\times$, at level $\Uone{Mp^s}$, $(\langle d\rangle f)(x)=\delta_{d,p}\cdot f(x\delta_d)$,
with $\delta_d\in\Uzero{Mp^s}$ having lower right entries $\equiv d$ (usually $\delta_{d,\ell}=\mathrm{diag}(1,d)$ at $\ell\mid Mp$, $1$ elsewhere;
the seed's ``right translation by $\mathrm{diag}(1,d)$'');
$\Pi_\ell:=T_{\Pi_\ell}$ for $\ell\mid D$, $(\Pi_\ell f)(x)=f(x\Pi_\ell)$;
and $U_p:=T_{\varpi_p}$, $\varpi_p=\iota_p^{-1}\mathrm{diag}(p,1)$, at levels $U^{(p)}K_s$ or $U^{(p)}K^0_s$ ($K_s=U_1(p^s)_p$, $K^0_s=U_0(p^s)_p$, $s\ge1$),
the normalisation $N_{U_p}$ of Def.~\ref{def:s01-NUp}.
For $\ell\nmid$ level, $[U\varpi_\ell U]=[U\,\mathrm{diag}(1,\ell)\,U]$; for $U_\ell$, $U_p$ the two choices differ (Lemma~\ref{lem:s01-Up}(c)).
\stproved{} (well-definedness). \stformal{} for trivial coefficients only (\lean{Langlandscape.Brandt.heckeT}, \lean{Langlandscape.Brandt.Tl},
\lean{Langlandscape.Brandt.Sl}, \lean{Langlandscape.Brandt.Ul}, \lean{Langlandscape.Brandt.diamond}, \lean{Langlandscape.Brandt.PiEll},
\lean{Langlandscape.Brandt.Up}).

\medskip
\noindent\rptnum{4.4}\label{rpt:prop:s01-relations}\textbf{Proposition 4.4 (relations).}
\emph{Let $U=\Uone{Mp^s}$ or $\Uzero{M}$, $V=V_k$.
\begin{enumerate}[label=(\alph*)]
\item $S_\ell=\ell^{k-2}\langle\ell\rangle^{-1}$ on $\Brandt{\Uone{Mp^s}}{V_k}$ ($\ell\nmid DMp$); $S_\ell=\ell^{k-2}$ at level
  $\Uzero{M}\cdot\GL_2(\Zp)$.
\item $\Pi_\ell^2=S_\ell$ (same formula as (a)) for $\ell\mid D$; $\Pi_\ell$ is independent of the choice of $\Pi_\ell$.
\item All operators of Def.~\ref{def:s01-hecke} commute; those supported at $\ell\ne q$ commute with the maps
  right translations $R_g:\Brandt{g^{-1}Ug}{V}\to\Brandt{U}{V}$, $f\mapsto f(\cdot\,g)$, and with $\rho_g:[x]\mapsto[xg]$,
  $\XB{U}\to\XB{g^{-1}Ug}$, for $g\in B_q^\times$, $q\ne p$: $T_h\circ R_g=R_g\circ T_h$ and $\rho_g\circ T_h=T_h\circ\rho_g$, with $T_h$
  taken at the respective levels. (So $R_gf=f\circ\rho_g=\rho_g^*f$ is pullback along $\rho_g$; this is the Lean convention.) For $g\notin N(\GL_2(\Z_q))$ the level
  $g^{-1}Ug$ is not contained in $\hat\cO^\times$; here Def.~\ref{def:s01-classset} and $\Gamma_x$ are used verbatim for the compact open subgroup $g^{-1}Ug\subseteq\Bhat^\times$
  (finiteness: $g^{-1}Ug\cap\hat\cO^\times$ has finite index in it), which is all that (c) needs.
\item $T_\ell$ for $\ell\nmid DMp$: $T_\ell T_{\ell^r}=T_{\ell^{r+1}}+\ell S_\ell T_{\ell^{r-1}}$ ($r\ge1$; $T_1=1$), with $T_{\ell^r}$ the sum of
  the double cosets $[U\gamma U]$, $\gamma$ supported at $\ell$, $\gamma_\ell\in M_2(\Z_\ell)$, $\det\gamma_\ell\in\ell^r\Z_\ell^\times$; with
  $T_n:=\prod_{\ell^r\Vert n}T_{\ell^r}$ for $(n,DMp)=1$, $T_mT_n=T_{mn}$ for coprime $m,n$; on $\Z^X$ at level $\Uzero{M}$ with $k=2$,
  $T_\ell\mathbf 1=(\ell+1)\mathbf 1$.
\end{enumerate}}
Status: (a)--(d) \stproved{}. At level $\hat\cO^\times$, (a) is \stformal{}; (d) is \stknown{}
\cite[Prop.~41.3.1(b), p.~770; Prop.~41.3.6 and (41.3.8), p.~771]{Vo} and \stformal{} for $(n,DM)=1$ at level $\hat\cO(M)^\times$. Also
\stformal{}: the second half of (b) (\lean{Langlandscape.Brandt.PiEll\_indep}), the first half of (c) for operators supported at two distinct
primes (\lean{Langlandscape.Brandt.heckeT\_comm\_of\_disjoint}), and the $R_g$-statement of (c) (\lean{Langlandscape.Brandt.heckeT\_comm\_rightTranslate}),
not the $\rho_g$-statement on $\XB{U}$.
\begin{outline}
(a) $z_\ell$ normalises $U$ and $z_{\ell,p}=1$, so $(S_\ell f)(x)=f(xz_\ell)=f(xy)$ with the scalar idele $y=\ell^{-1}z_\ell$ (the global
$\ell^{-1}$ is central); $y=\delta_{\ell^{-1}}u$ with $u\in\Uone{Mp^s}$, and $y_p^{-1}=\ell I$ acts on degree-$n$ polynomials by $\ell^{n}$,
giving $S_\ell=\ell^{k-2}\langle\ell\rangle^{-1}$; at level $\Uzero{M}\cdot\GL_2(\Zp)$ one has $y\in U$. (b) $\Pi_\ell^2=z_\ell v$ with
$v\in\cO_\ell^\times\subseteq U$, $v_p=1$; two uniformisers differ by a unit. (c) Two principles: operators supported at disjoint sets of
primes commute (choose disjointly supported representatives), and $T_\delta T_g=T_{\delta g}=T_{g\delta}=T_gT_\delta$ for $\delta$
normalising $U$ and commuting with $g$; for $R_g$, the representatives $h_i$ (supported at $\ell\ne q$) commute with $g$ and represent the
operator at both levels. (d) Locally, $\Delta_r/\GL_2(\Z_\ell)$ is the set of lattices of index $\ell^r$ in $\Z_\ell^2$; counting the
chains $L'\subseteq L_1\subseteq\Z_\ell^2$ ($1$ or $\ell+1$ choices of $L_1$) gives the recursion, and $T_\ell\mathbf 1=\#\{\text{index-}\ell\text{ lattices}\}=\ell+1$.
\end{outline}

\noindent\rptnum{4.5}\label{rpt:rem:s01-diamond}\textbf{Remark 4.5 (diamond convention flag).}
With the seed's adelic diamond (right translation by $\mathrm{diag}(1,d)$), Prop.~\ref{prop:s01-relations}(a) reads
$S_\ell=\ell^{k-2}\langle\ell\rangle^{-1}$ (\stproved), whereas classically $S_\ell=\ell^{k-2}\langle\ell\rangle$ on $S_k(\Gamma_1(N))$
(\stknown; compare \cite[\S4, p.~13]{Buz04}; the precise comparison is Lemma 5.5(a) of \mainpdf). Since the Jacquet--Langlands
correspondence of Thm.~\ref{thm:s01-JLclassical} is compatible with $T_\ell$ and $S_\ell$, the adelic $\langle d\rangle$ corresponds to the
classical $\langle d\rangle^{-1}$ (\stpmc{} from the two formulas and the cited compatibility). The seed is not changed; stages 04 and 05
use $\langle d\rangle^{-1}$ when comparing $\Lambda$-structures with $\GL_2$. The same inversion reappears in Lemma~\ref{lem:s01-Up}(d):
$U_p^*=\langle p\rangle_MU'_p$ involves the \emph{adelic} $\langle p\rangle_M$.

\medskip
\noindent\rptnum{4.6}\label{rpt:def:s01-T}\textbf{Definition 4.6 (Hecke algebras)} (condensed).
Let $V$ be $\Zp$-torsion free (trivial or $V_k$). $\T_B(U;V)\subseteq\End\Brandt{U}{V}$ is the $\Zp$-algebra generated by
$T_\ell,S_\ell^{\pm1}$ ($\ell\nmid DMpq$), $U_\ell$ ($\ell\mid M$), $\langle d\rangle$, $U_p$ (if $s\ge1$) and $\Pi_\ell$ ($\ell\mid D$); the anemic
version $\T_B^{\mathrm{an}}$ omits $U_\ell,U_p,\Pi_\ell$. It is commutative (Prop.~\ref{prop:s01-relations}(c)) and finite free over $\Zp$.
\emph{Relation with $\nu$} (trivial coefficients, $p\nmid w_x$ for all $x$): $T\mapsto\nu T\nu^{-1}$ is a ring isomorphism from $\T_B$ on
$\Zp[X]$ onto the algebra $\T_B^\dagger$ on $\Zp^X$ generated by the $T_{g^{-1}}$: $T_\ell\mapsto S_\ell^{-1}T_\ell$, $S_\ell\mapsto S_\ell^{-1}$,
$\langle d\rangle\mapsto\langle d\rangle^{-1}$, $\Pi_\ell\mapsto\Pi_\ell^{-1}$, $U_p\mapsto\langle p\rangle_MU'_p$ (Lemma~\ref{lem:s01-Up}(d)),
$U_\ell\mapsto\langle\ell\rangle_{(\ell)}U'_\ell$ for $\ell\mid M$, with $U'_\ell=T_{\eta_\ell}$, $\eta_\ell=\iota_\ell^{-1}\mathrm{diag}(1,\ell)$. So the
anemic algebra is stable, but $U_p$, $U_\ell$ go to the primed operators, which are in general \emph{not} in $\T_B$: $\nu$ identifies $\Zp[X]$
and $\Zp^X$ as $\T_B$-modules only after twisting by $T\mapsto T^\dagger$. \stproved{}. No Lean counterpart in stage 01.

\subsubsection*{Chapter 2: the tree and the $S$-arithmetic group}

\noindent\rptnum{7.1}\label{rpt:def:s02-tree}\textbf{Definition 7.1 (Bruhat--Tits tree, types, base cells).}
$\cT_q$ is the graph whose vertices are the homothety classes $[L]$ of $\Zq$-lattices $L\subset\Qq^2$ (of rank~$2$).
$[L]\ne[L']$ are adjacent iff there are representatives with $qL\subsetneq L'\subsetneq L$.
$\GL_2(\Qq)$ acts by $g[L]=[gL]$. Every lattice is of the form $g\Zq^2$ with $g\in\GL_2(\Qq)$, and $g\Zq^2=g'\Zq^2$ iff
$g'\in g\GL_2(\Zq)$. The \emph{type} is $\tau[g\Zq^2]:=\ord_q\det g\bmod 2$ (well defined).
Put $v_0=[\Zq^2]$, $v_1=\alpha v_0=[\Zq\oplus q\Zq]$ with $\alpha=\mathrm{diag}(1,q)$, and $e_0=\{v_0,v_1\}$. Every edge has one vertex of
each type (Lemma~\ref{lem:s02-tree}). Every edge is oriented from its type-$0$ end $o(e)$ to its type-$1$ end $t(e)$.

\medskip
\noindent\rptnum{7.4}\label{rpt:def:s02-GU}\textbf{Definition 7.4 (admissible levels; $\GU$)} (condensed).
A level $U\subseteq\hat\cO^\times$ is \emph{$q$-admissible} if it is compact open, $U=U_q\times U^{(q)}$ with
$U_q=\iota_q^{-1}\GL_2(\Zq)$, and $\Nrd(U)=\Zhat^\times$. All levels $\Uzero{N},\Uone{N}$ with $q\nmid N$ and
$U_s=\Uone{Mp^s}$ are $q$-admissible (Lemma 1.3 of \mainpdf; $q\nmid Mp$). For such $U$ put
\[\GU:=B^1\cap\big(B^1_q\times U^{(q)}\big)=\{\gamma\in B^\times:\Nrd\gamma=1,\ \gamma^{(q)}\in U^{(q)}\}=\cO_U[\tfrac1q]^1\]
(the last expression is notation; no order $\cO_U$ is asserted for general $U$). For $U=\hat\cO^\times$ it is the group of norm-one units
of $\cO[\frac1q]$. $\GU$ acts on $\cT_q$ through $\iota_q$, type-preservingly and without inversions, and on a $\Zp[U_p]$-module $V$
through $\gamma\mapsto\gamma_p$. The levels attached to the cells are $U$ (type $0$), $\alpha U\alpha^{-1}$ (type $1$) and
$\Uzero{q}U:=U\cap\Uzero{q}$ (edges), all with $\Nrd=\Zhat^\times$. \stproved{}. \stformal{} (definitions and the level statements,
\lean{Langlandscape.Tree.GammaU}, \lean{Langlandscape.Tree.IsQAdmissible} and relatives).

\medskip
\noindent\rptnum{7.8}\label{rpt:def:s02-BG}\textbf{Definition 7.8 (bipartite Brandt graph).}
For $q$-admissible $U$, $\BG(U,q)$ is the finite bipartite multigraph with $V_0=V_1=\XB{U}$ (we write $(x,0)$ and $(x,1)$ for
$x\in\XB U$), edge set $E=\XB{U'}$, $U'=\Uzero{q}U$, and
\[o[y]=(s_0[y],0),\quad t[y]=(s_1[y],1),\qquad s_0[y]=[y],\qquad s_1[y]=[y\alpha]\qquad(\XB{U'}\to\XB{U}).\]
The weights are $w_{(x,j)}=w_x$ (level $U$) and $w_{[y]}=w_y$ (level $U'$), with $w_x=|\Gamma_x/(\Gamma_x\cap\{\pm1\})|$ as in
Def.~\ref{def:s01-classset}. This is the graph that HACS (and the Python track) computes.
\stproved{} (well-definedness and finiteness); \stformal{} (graph, well-definedness, finiteness; \lean{Langlandscape.Tree.BG}; the weights
are not part of the Lean structure).

\medskip
\noindent\rptnum{7.23}\label{rpt:rem:s02-torsion}\textbf{Remark 7.23 (torsion in $H_1(\GU,\Z)$; deviation D1).}
By Prop.~\ref{prop:s02-coh}(a), $H_1(\GU,\Z)$ is an extension of the free group $H_1(\BG,\Z)$ by the cokernel of
$\bigoplus_eH_1(\Gamma_e,\Z)\to\bigoplus_vH_1(\Gamma_v,\Z)$, which is finite and killed by $24$, and non-zero in general. \stcomp{}
Example: $D=2$, Hurwitz order, $U=\hat\cO^\times$, $q=3$: one vertex orbit per type with $\Gamma_v\cong\SL_2(\F_3)$, one edge orbit with
$\Gamma_e$ cyclic of order $6$, so $\GU\cong\SL_2(\F_3)*_{C_6}\SL_2(\F_3)$, $\BG(U,3)$ is a single edge ($H_1(\BG,\Z)=0$, matching genus $0$ of
the Shimura curve of discriminant $6$), and $H_1(\GU,\Z)\cong\Z/3$ (result file \res{s02/torsion.json}).
\emph{D1:} since $H_1(\GU,\Z)$ has torsion, it is not a character group; over $\Z$ the indefinite stage is taken to be
$\cX^{(q)}_{Dq}(U):=H_1(\BG(U,q),\Z)=H_1(\GU,\Z)/\mathrm{torsion}$ (Def.~7.21 of \mainpdf), working with the two-term complex of
Prop.~\ref{prop:s02-coh}. This agrees with $H_1(\GU,\Z)$ after $\otimes\Zp$ for $p\ge5$.

\medskip
\noindent\rptnum{7.24}\label{rpt:rem:s02-CD}\textbf{Remark 7.24 (\v{C}erednik--Drinfeld; cited, never used).}
If $\Nrd(U^{(q)})$ is maximal and $U^{(q)}$ contains rational scalars of every $q$-adic valuation ($\Gamma_0$-type levels qualify,
$\Uone{N}$-levels do not), the \v{C}erednik--Drinfeld uniformisation \cite[Ch.~III, Th\'eor\`eme~5.2, pp.~140--141]{BC} describes the
Shimura curve of discriminant $Dq$ at $q$; it is twisted (Frobenius swaps the two vertex types). The dual graph of the special fibre is
$\BG(U,q)$ with edge lengths $|\Gamma_x/\{\pm1\}|$ \cite[Ch.~III, 5.3.3, p.~143; 5.4.3--5.4.4, pp.~145--146]{BC}; the description of the
character group $\cX^{(q)}_{Dq}(U)$ as $H_1$ of this graph is derived later in \cite{BC} (Ribet's chapter, p.~159, via SGA~7). \stcited{}
The tree route never uses this remark.

\subsubsection*{Chapter 3: graphs, Ribet, Ihara}

\noindent\rptnum{8.1}\label{rpt:def:s03-graph}\textbf{Definition 8.1 (weighted graph).}
A \emph{weighted graph} $\G$ is a finite graph (multiple edges and loops allowed) with vertex set $V(\G)$, edge set $E(\G)$, a
chosen orientation, and weights $w_e\in\Z_{\ge1}$. Its chain complex is $\Z[E]\xrightarrow\partial\Z[V]$, $\partial e=t(e)-o(e)$. Put
$X(\G):=H_1(\G)=\ker\partial$, $\langle e,e'\rangle_w=w_e\delta_{ee'}$ and
\[\Phicomb(\G):=\coker\big(X(\G)\to\Hom(X(\G),\Z),\ x\mapsto\langle x,\cdot\rangle_w\big).\]
\status{definition} \stformal{} (definitions; \lean{Langlandscape.Graph.WGraph}, \lean{Langlandscape.Graph.phiComb}).
Basic facts (\stproved): $X(\G)$ is a free direct summand of $\Z[E]$; $|\Phicomb(\G)|=\det G>0$ for the Gram matrix $G$ of a basis of
$X(\G)$; $\Phicomb(\G)$ does not depend on the orientation up to isomorphism.

\medskip
\noindent\rptnum{8.2}\label{rpt:rem:s03-grothendieck}\textbf{Remark 8.2 (graphs and semistable curves).}
For a semistable curve with dual graph $\G$ and thicknesses $w_e$, $X(\G)$ is the character group of the toric part, $\langle,\rangle_w$ is
the monodromy pairing and $\Phicomb$ is the component group: \cite[IX]{SGA7} \stcited{}, and \cite[Prop.~2.2, Thm.~2.3, p.~436; weighted
form pp.~437--438]{Ri90}. The tree route does not use this dictionary; in the plan it is a remark only, and its use belongs to stage 05.

\medskip
\noindent\rptnum{8.6}\label{rpt:def:s03-DR}\textbf{Definition 8.6 (DR graph).}
Let $B=B_q$ and $U=\Uone{Mp^s}$ or $\Uzero M$. Then $\GDR_q(U)$ is the two-vertex weighted graph with edge set $\XBq U$, every edge
oriented $0\to1$, and $w_e=|\Gamma_e/(\Gamma_e\cap\{\pm1\})|$. Put $\Xcomb_q(U):=\Z[\XBq U]^0$ and, for $k>2$,
$\Xcomb_q(U;\Vk):=\Brandt U{\Vk}$. \status{definition} \stformal{} at $U=\hat\cO^\times$ (\lean{Langlandscape.Ribet.DRgraph}).
(For $s\ge1$ every $w_e=1$, so $\Phicomb(\GDR_q(U))\cong\Z/|S_s|$ by Prop.~\ref{prop:s03-phi}, and $X(\GDR_q(U))=\Xcomb_q(U)$; \stproved.)

\medskip
\noindent\rptnum{12}\label{rpt:ss:s03-amalgam}\textbf{Section 12 (the $\GL_2$ side: the amalgam, Betti only).}
For $(M,q)=1$, $\Gamma=\Gamma_0(M)$, $G=\Gamma_0(M)[\frac1q]$, $\alpha'=\mathrm{diag}(q,1)$ and $\Gamma'=\alpha'^{-1}\Gamma\alpha'$, this chapter
section proves the amalgam $G=\Gamma*_{\Gamma_0(Mq)}\Gamma'$ (Prop.~\ref{prop:s03-amalgam}), the Mayer--Vietoris sequence with
$V=\Sym^{k-2}$, $k\ge4$ even, $l>k$, $l\ne q$ (Prop.~\ref{prop:s03-MV}), $H^1(G,V)=0$ (Lemma 12.8 of \mainpdf), the formula
$\Psi^*\Psi=\left(\begin{smallmatrix}q+1&-q^{2-k}T_q\\-T_q&q+1\end{smallmatrix}\right)$ (Prop.~\ref{prop:s03-PsiPsi}) and the Hecke
equivariance of $\Psi$, all \stpmc; it defines the Betti level-raising module $\mathrm{LR}_q(M,k)_\m=\coker(\Psi^*\Psi)_\m$ (finite, \stproved)
and states the Betti Ihara lemma for $\GL_2$ as \stexp{} (Conj.~\ref{conj:s03-iharaGL2}). The identification of $T_{\alpha'}$ with the
classical $T_q$ is \stpmc{} and \stcomp{} on the HACS cases.

\medskip
\noindent\rptnum{12.2}\label{rpt:prop:s03-amalgam}\textbf{Proposition 12.2 (amalgam).}
\emph{$G=\Gamma*_{\Gamma_0(Mq)}\Gamma'$.} \stpmc: \cite[II.1.4, Thm.~3, p.~79]{Se} under hypothesis (D) of \cite[II.1.4, Thm.~2, p.~78]{Se}:
the closure of $G$ in $\SL_2(\Qq)$ is $\SL_2(\Qq)$ (Lemma 12.1(d) of \mainpdf, proved elementarily; alternatively from strong approximation
for $\SL_2$ over $\Q$, \cite[Thm.~28.2.3, p.~480]{Vo}). For $M=1$ it is \cite[II.1.4, Cor.~2, p.~80; p.~81]{Se} \stcited, whose second
factor is exactly $\{(\begin{smallmatrix}a&q^{-1}b\\qc&d\end{smallmatrix})\}=\Gamma'$. The vertex stabilisers are $\Gamma$ (of $v_0$) and
$\Gamma'$ (of $v_1$).
\begin{outline}
$G\subseteq\SL_2(\Qq)$ acts on $\cT_q$ without inversion (Lemma~\ref{lem:s02-tree}(b)). Serre's Thm.~3 says that under (D), $G$ is the amalgam
of the stabilisers of $v_0$ and $v_1$ along the stabiliser of the edge. These are: $G\cap\SL_2(\Zq)=G\cap\SL_2(\Z)=\Gamma$;
$G\cap\Stab(v_1)=\alpha'^{-1}(\alpha'G\alpha'^{-1}\cap\SL_2(\Zq))\alpha'=\alpha'^{-1}\Gamma\alpha'=\Gamma'$ (using $\alpha'G\alpha'^{-1}=G$);
and the edge stabiliser $\Gamma\cap\Gamma'=\Gamma_0(Mq)$. Hypothesis (D) holds because $G$ contains the elementary matrices with entries in
the dense subsets $\Z[\frac1q]$, $M\Z[\frac1q]$ of $\Qq$, which generate $\SL_2(\Qq)$ topologically.
\end{outline}

\noindent\rptnum{12.3}\label{rpt:prop:s03-MV}\textbf{Proposition 12.3 (Mayer--Vietoris).}
(Hypotheses: $k\ge4$ even, $l$ prime, $l>k$, $l\ne q$, $(M,q)=1$, $V=\Sym^{k-2}(\Z_{(l)}^2)$ the standard column representation,
$c:H^1(\Gamma,V)\xrightarrow\sim H^1(\Gamma',V)$ the conjugation isomorphism.)
\emph{There is an exact sequence, with $\Psi(x,y)=\mathrm{res}\,x-\mathrm{res}\,(cy)$,
\[0\to H^1(G,V)\to H^1(\Gamma,V)^2\xrightarrow{\ \Psi\ }H^1(\Gamma_0(Mq),V)\to H^2(G,V)\to0.\]}
\stpmc: the Mayer--Vietoris sequence of an amalgam with coefficients \cite[III.6, Exercise, p.~74]{Br} \stcited{} (the trivial-coefficient
version is \cite[II, Cor.~7.7, p.~51]{Br}), applied to Prop.~\ref{prop:s03-amalgam}. The end terms vanish for these reasons.
$H^0(\Gamma_0(Mq),V)=0$ for $k>2$: the image of $\Gamma_0(Mq)$ in $\SL_2(\F_l)$ is everything if $l\nmid M$, and $\Sym^{k-2}\F_l^2$ is irreducible and
non-trivial since $k-2<l$; if $l\mid M$ the image lies in the Borel subgroup, and the invariants vanish as well (proof). $H^2(\Gamma,V)=0$ because
$\Gamma$ is virtually free with torsion of order prime to $l$. The same argument shows that $H^1(\Gamma,V)$ is torsion-free.
Moreover $H^1(G,V)=0$ (Lemma 12.8 of \mainpdf), so the sequence reads
\[0\to H^1(\Gamma,V)^2\xrightarrow{\ \Psi\ }H^1(\Gamma_0(Mq),V)\to H^2(G,V)\to0.\]
\begin{outline}
The Mayer--Vietoris sequence of $G=\Gamma*_{\Gamma_0(Mq)}\Gamma'$ has middle map a difference of restrictions; identifying $H^1(\Gamma',V)$
with $H^1(\Gamma,V)$ via $c$ gives $\Psi$. $H^0(\Gamma_0(Mq),V)=0$: for $l\nmid M$ the reduction map to $\SL_2(\F_l)$ is onto and
$\Sym^n\F_l^2$ ($n=k-2<l$) is irreducible and non-trivial; for $l\mid M$ the unipotent element fixes only the line $\F_le_1^n$, on which a
suitable $\gamma\in\Gamma_0(Mq)$ with upper left entry a generator of $\F_l^\times$ acts by $a^n\ne1$; a non-zero integral invariant would
reduce to a non-zero invariant mod $l$. $H^i(\Gamma,V)=0$ for $i\ge2$: $\Gamma$ acts on the tree of $\SL_2(\Z)$ with finite stabilisers of
order prime to $l$ \cite[I.4.2(c), p.~35]{Se}. Torsion-freeness of $H^1(\Gamma,V)$ follows from $H^0(\Gamma,V/l)=0$.
\end{outline}

\subsubsection*{Chapter 4: Hida theory on the definite side}

\noindent\rptnum{13.1}\label{rpt:ss:s04-notation}\textbf{Section 13.1 (standing notation of stage 04).}
It fixes $p\ge5$, $p\nmid D$, $\cO=\cOE$ the integers of a finite $E/\Qp$ with uniformiser $\varpi$, $\Lam=\cOE[[1+p\Zp]]$, levels prime
to $p$ $U=U^{(p)}\times\GL_2(\Zp)$ with $\Nrd U=\Zhat^\times$, the tower $UU_s=U^{(p)}\times\Kp{s}$ ($\Kp{s}=\iota_p^{-1}\Uone{p^s}_p$),
$X_s(U)=\XB{UU_s}$, $X'_s(U)=\XB{U^{(p)}\Kpp{s}}$, and the projection $\pi:X_{s+1}\to X_s$. It records (all from stage 01, \stproved{}
there) that for $s\ge1$ all $\Gamma_x$ are trivial ($w_x=1$, $\nu=$ identity on bases), $\Delta_s$ acts freely by adelic diamonds,
$|X_{s+1}|=p^2|X_s|$, and $U_p=\NUp$ commutes with $\pi_*$, the diamonds, $T_\ell$, $S_\ell$, $\Pi_\ell$ and right translations away from
$p$. Modules are homological unless marked otherwise.

\medskip
\noindent\rptnum{13.1}\label{rpt:def:s04-e}\textbf{Definition 13.1 (ordinary projector).}
For a finitely generated $\Zp$-module $M$ with an endomorphism $T$, $e_T:=\lim_{n\to\infty}T^{n!}$. It exists, is an idempotent
in the closed subalgebra $\Zp[[T]]\subseteq\End M$ generated by $T$, and $M=e_TM\oplus(1-e_T)M$ with $T$ invertible on $e_TM$ and
topologically nilpotent on $(1-e_T)M$. The functor $M\mapsto e_TM$ is exact on $\Zp[T]$-modules of this kind. If $M$ is finite,
$e_T=T^{N}$ for any $N$ that is a multiple of the period and at least the index of $T$ in the finite monoid $\End M$
(e.g.\ $N=|\End M|!$). We write $e:=e_{U_p}$ and $e^*:=e_{U_p^*}$, where $U_p^*$ is the $\langle,\rangle$-adjoint of $U_p$; at level
$\Uone{Mp^s}$, $U_p^*=\langle p\rangle_MU'_p$ (Lemma~\ref{lem:s01-Up}(d)), and Lemma~\ref{lem:s04-ekills}(d) gives the same formula at
every level $UU_s$.
\stproved{}; the existence of the limit is also \stknown{} \cite[\S7.2, Lemma~1, p.~201]{Hi93}. \stformal{} for $\cO=\Zp$ in the form of a
decomposition $M=\mathrm{ord}\oplus\mathrm{nil}$ with projector (\lean{Langlandscape.Ordinary.ordProj}); the limit formula itself and
general $\cO_E$ are not formalized.

\medskip
\noindent\rptnum{14.1}\label{rpt:def:s04-Mord}\textbf{Definition 14.1 ($\Lambda$-adic Brandt module).}
For $U$ of level prime to $p$: $\Mord(U):=\varprojlim_{s\ge1}e\,\cOE[X_s(U)]$, limit along $\pi_*$ (which commutes with $U_p$ by
Lemma~\ref{lem:s01-Up}(b), and with all $T_\ell,S_\ell,U_\ell,\Pi_\ell,\langle d\rangle$). $\Lam$ acts through the adelic diamonds:
$[u]\cdot a:=\langle u\rangle a$, $u\in1+p\Zp$. The finite group $\mu_{p-1}$ acts as well, and $\Mord(U)$ decomposes integrally into
$\mu_{p-1}$-eigenpieces $\Mord(U)^{(\chi)}$. For $U=\Uone M$ also $(\Z/M)^\times$ acts; if $p\nmid\varphi(M)$, $\Mord(U)$ decomposes
integrally into eigenpieces for $\mu_{p-1}\times(\Z/M)^\times$, and if $p\mid\varphi(M)$ only after inverting $p$ (as in
Thm.~\ref{thm:s04-control}(b)).
Well-definedness \stproved. \stformal{} for $E=\Qp$ and $U=\Uone M$, as compatible families (\lean{Langlandscape.Ordinary.Mord},
\lean{Langlandscape.Ordinary.Mord.lamModule}), with the commutation of the diamonds with $\pi_*$ as an explicit argument; general $\cO_E$
and the eigenpiece remark are not formalized.

\medskip
\noindent\rptnum{16.1}\label{rpt:def:s04-Pk}\textbf{Definition 16.1 (arithmetic points in the seed's conventions).}
For $k\ge2$ put $P^{\rm ad}_k:=\ker\big(\Lam\to\cOE,\ [u]\mapsto u^{-(k-2)}\big)$. The seed's $P_k=\ker([u]\mapsto u^{k-2})$ is the
image of $P^{\rm ad}_k$ under the involution $\iota:[u]\mapsto[u^{-1}]$, i.e.\ it is the right ideal for the $\Lam$-structure via
\emph{classical} diamonds $\langle u\rangle_{\rm cl}=\langle u\rangle^{-1}$.
Justification \stproved{}: via $\mathrm{ev}$ of Equation~(1), the adelic diamond $\langle d\rangle$, $d\equiv1\bmod p$, acts by
$d^{-(k-2)}$ on functions coming from weight $k$. This matches \cite[\S4.4, Def.~4.1, p.~31]{Hs}. The weight-$k$ point with trivial
character at $p$ also requires $\mu_{p-1}$ to act by $\omega^{-(k-2)}$.

\subsubsection*{Chapter 5: the comparison hypothesis}

\noindent\rptnum{18.2}\label{rpt:ss:s05-sign}\textbf{Section 18.2 (setting, the dictionary, and the sign).}
It fixes the setting of stage 05 ($B=B_q$, $U_s=\Uone{Mp^s}$, $S_s=\XBq{U_s}$; on the $\GL_2$ side $\Gamma_s=\Gamma_1(Mp^s)\cap\Gamma_0(q)$,
$J_s=\Jac X_s$ in the $\mu_N$-model with Picard functoriality, the model being a team decision (FD1), and the character group
$\Xq(J_s)$ at $q$). It proves the factor-$2$ bookkeeping between $\Gamma_0$- and $\Gamma_1$-type levels (\stproved; Hypothesis (F2)),
defines the dictionary $\theta_s$ (Def.~\ref{def:s05-theta}) and proves it well defined (\stproved), and proves the sign
$\varepsilon=+1$ and the exponent statements of Prop.~\ref{prop:s05-sign} (\stpmc{}, the $\Gamma_1$-level Frobenius part given (X1) and
FD1). It also records the status of $\HypX$ with the four-row Frobenius table (Remark 18.10 of \mainpdf; $\HypX$ stays \sthyp) and an
internal consistency lemma for $\theta_s$.

\medskip
\noindent\rptnum{18.4}\label{rpt:def:s05-theta}\textbf{Definition 18.4 (dictionary $\theta_s$).}
Let $\cH_s:=\Z[(\Z/Mp^s)^\times]\otimes_\Z\Z[T_\ell,U_\ell,U_p\text{ or }T_p,U_q]$, the group ring of the diamonds $\langle d\rangle$
($d\in(\Z/Mp^s)^\times$) tensored with the polynomial ring on the symbols
$T_\ell$ ($\ell\nmid Mpq$ prime), $U_\ell$ ($\ell\mid M$), $U_p$ (if $s\ge1$), $T_p$ (if $s=0$), $U_q$.
\revnote{S05-M1: the stage tex takes the free commutative algebra on all symbols including the $\langle d\rangle$; then
$\langle d\rangle^{-1}$, $T_\ell^\Ros=\langle\ell\rangle^{-1}T_\ell$ and $U_q^\Ros$ would not lie in $\cH_s$. The group ring fixes this; $\theta_s$ is well defined on it
because $d\mapsto\langle d\rangle^{-1}$ (adelic) is a homomorphism into a commutative group (self-audit F11).}
It acts on $\Xq(J_s)$ through the classical Hecke correspondences (Picard functoriality, classical normalisation, classical
diamonds). Define $\theta_s:\cH_s\to\End_\Z\Z[S_s]^0$ as the $\Z$-algebra homomorphism with
\[T_\ell\mapsto T_\ell,\quad U_\ell\mapsto U_\ell,\quad U_p\mapsto U_p\ (N_{U_p}),\quad T_p\mapsto T_p,\quad
\langle d\rangle\mapsto\langle d\rangle^{-1}\ (\text{adelic}),\quad U_q\mapsto\varepsilon\,\Pi_q,\quad \varepsilon=+1,\]
all acting homologically on $\Z[S_s]$ and restricted to $\Z[S_s]^0$ (operators as in Def.~\ref{def:s01-hecke}; at $s=0$,
$T_p:=T_{\varpi_p}$ at the level $U_0$, whose $p$-component is $\GL_2(\Zp)$). For $M=1$ there are no $U_\ell$; for $s=0$ and $M\le2$ the
diamond group is trivial. \stformal{} on the \emph{generators} only (\lean{Langlandscape.Comparison.theta}; no ring-homomorphism extension
in Lean, without $U_\ell$ for $\ell\mid M$).

\medskip
\noindent\rptnum{18.5.1}\label{rpt:sss:s05-general}\textbf{Section 18.5.1 (several ramified primes).}
It shows (Prop.~\ref{prop:s05-nochain}) that tree quotients and Ribet sequences only ever produce subquotients of Brandt modules of the
\emph{same} definite algebra (\stproved), that this is no obstruction at the level of isomorphism classes, since for $D=q_1q_2q_3$ classical
JL applied twice gives a rational Hecke isomorphism (\stpmc), and hence that the true content is the absence of a \emph{canonical}
comparison map. It then tabulates the geometric chain for $D=q_1q_2q_3$, whose links need $\HypX$, \v{C}erednik--Drinfeld and Galois
input; the several-primes defect formula is stage 06's (\stconj{}/\stsketch{} there).

\medskip
\noindent\rptnum{18.7}\label{rpt:ss:s05-status}\textbf{Section 18.7 (what the literature states at $\Gamma_1$-level).}
It splits $\HypX$ at level $\Gamma_1(Mp^s)\cap\Gamma_0(q)$ into four inputs and records the literature found: (S1) the
Deligne--Rapoport semistable model is \stcited{} \cite[Ch.~VI, Thm.~6.9]{DR}; (S2) Deuring at $\Gamma_1$-level with $U_p$ and diamonds is
stated nowhere looked at; (S3) the monodromy pairing follows from (S1)+(S2) by Raynaud's theorem (\stpmc, our inference); (S4) the Frobenius
clause is \stpmc{} (the four-row table, derived from the naive/Albanese row, with \cite[Lemma~3.17]{RS} only as a consistency check),
conditional on FD1. (S2) and the naturality in $s$ (X3) remain \stexp, so $\HypX$ stays a hypothesis.

\medskip
\noindent\rptnum{18.38}\label{rpt:rem:s05-degree}\textbf{Remark 18.38 (no Betti-canonical replacement).}
After inverting $q$, the cuspidal cohomology of $\GL_2$ sits in degree $2$ (the $q$-new part of $H^*(\Gamma_0(M)[\frac1q],V)$ is in $H^2$,
Section~12) and that of the definite $S$-arithmetic group in degree $1$ (Prop.~\ref{prop:s02-coh}); for $B_q$ itself there is no tree and
it sits in degree $0$ (Prop.~\ref{prop:s05-nochain}(a)). So a degree-preserving Hecke-equivariant chain map inducing $\HypX$ on cuspidal
parts cannot exist between the $S$-arithmetic complexes of $\GL_2(\Z[\frac1q])$-type and of $B_q$; any realisation must shift degree
(e.g.\ through a boundary or Gysin-type map) or go through a third object. \stsketch{} (degree count only).

\medskip
\noindent\rptnum{18.39}\label{rpt:def:s05-stages}\textbf{Definition 18.39 (the four stages as simplicial orbifolds).}
\begin{enumerate}[label=(\roman*)]
\item Definite level $U$: the $0$-dimensional orbifold (groupoid) $\coprod_{[x]\in\XB{U}}B\Gamma_x$; its $H^0$ with coefficients $V$ is
  $\Brandt{U}{V}$.
\item Definite $S$-arithmetic stage: the graph of finite groups $\Gamma_U\quot\cT_q$ (Prop.~\ref{prop:s02-quot}) with structure maps
  $s_0,s_1$; a $1$-dimensional orbifold model of $B\Gamma_U$ after barycentric subdivision, orbifolds being taken as étale groupoids
  \cite[\S4, Thm.~4.1, p.~11]{MP97} (for the simplicial model and its cohomology with local coefficients: Moerdijk--Pronk 1999, Thm.~2.1.1,
  p.~278). Over $\Z$ its $H_1$ is taken as in Remark 18.35 of \mainpdf{} (i.e.\ modulo torsion, as in Remark~\ref{rem:s02-torsion}).
\item $\GL_2$ arithmetic stage: $\Gamma_s\backslash\Hb$ with a finite triangulation (as in HACS).
\item $\GL_2$ $S$-arithmetic stage: the graph of spaces $[Y_0(M)\sqcup Y_0(M)\leftleftarrows Y_0(Mq)]$ (Section~12).
\end{enumerate}

\medskip
\noindent\rptnum{18.42}\label{rpt:rem:s05-O46}\textbf{Remark 18.42 (O4, O6 made precise).}
O4: a Shimura-curve $\HypX$ for discriminant $d$ and $r\nmid d$ is Hypothesis~\ref{hyp:s05-X} with $J_s$ replaced by the Jacobian of the
Shimura curve of discriminant $d$ and level $\Gamma_1(Mp^s)\cap\Gamma_0(r)$, $B=B_{dr}$; at $\Gamma_0$-level it is
\cite[Thm.~4.3, p.~11]{Koh01} (Buzzard). O6: at an Eisenstein $\m$, (X2) carries information that $\HypXiso$ does not; its weight-$2$,
level-$1$ shadow is the Mazur check $\Phicomb(\GDR_q(1))\cong\Z/\mathrm{num}\frac{q-1}{12}$ versus $\Phi_q(J_0(q))$
\cite[Appendix, Thm.~(A.1)(b), p.~173]{Mz}; this is a check of (X2) at $\Gamma_0$-level, not a proof.

\subsubsection*{Chapter 6: congruence ideals and the defect}

\noindent\rptnum{19.1}\label{rpt:def:s06-gor}\textbf{Definition 19.1 (traces, Gorenstein, congruence ideal).}
Let $A$ be a commutative ring, $S$ a commutative $A$-algebra, $S^\vee:=\Hom_A(S,A)$.
\begin{enumerate}[label=(\alph*)]
\item The dual $S^\vee$ is an $S$-module by $(s\varphi)(x):=\varphi(sx)$. An element $\ell\in S^\vee$ is a \emph{trace} if $S$ is finite
  free over $A$ and the $S$-linear map $S\to S^\vee$, $s\mapsto s\ell=\ell(s\,\cdot)$, is bijective. Then $S^\vee$ is free of rank one over
  $S$ with generator $\ell$. $S$ is \emph{$A$-Gorenstein} if it has a trace.
\item For an ideal $I\subseteq S$ put $\Ann_S(I):=\{s\in S:sI=0\}$.
\item For an $A$-algebra map $\mu:S\to A$ put $S[\mu]:=\Ann_S(\ker\mu)$ and $\eta(S,\mu):=\mu(S[\mu])\subseteq A$, the
  \emph{congruence ideal}.
\end{enumerate}
\status{definition} (Lean: \texttt{IsGorensteinWith}, \texttt{IsGorenstein}, \texttt{ann}, \texttt{congruenceIdeal} in
\pathtt{Langlandscape.CommAlg}.)

\medskip
\noindent\rptnum{20.1}\label{rpt:ss:s06-setting}\textbf{Section 20.1 (setting and standing hypotheses).}
It fixes $D=q$, $B=B_q$, $p\ge5$, $q\ne p$, $(M,pq)=1$, $N=qM$, $\Lam=\cOE[[1+p\Zp]]$ acting through the \emph{classical} diamonds on the
$\GL_2$ side (weight-$k$ point $P_k$; the definite side's adelic $P^{\rm ad}_k$ is transported by $[u]\mapsto[u]^{-1}$), and the element
$S_q\in\Lam^\times$ with $P_k(S_q)=q^{k-2}$. It states the standing hypotheses (H0), (H0$^-$), (Lev), (H1), (H2), (Im), (P) of
Hypothesis~\ref{hyp:s06-H} (\sthyp), and shows that (H2)$+$(H0) imply Hsieh's (CR,$\{q\}$), so that the combinatorial half needs only (H0),
(H1), (H2) (\stproved{} for the implication; the hypothesis itself is \stcited, \cite[\S1.4]{Hs}).

\medskip
\noindent\rptnum{20.3}\label{rpt:def:s06-rings}\textbf{Definition 20.3 (Hecke rings and congruence ideals).}
$\T'$ and $\T''$ are the images of $\T$ in the endomorphisms of the $q$-new and of the $q$-old ordinary $\Lam$-adic cusp forms
(localised at $\m$). Put $\fa:=\ker(\T\to\T')$ and $\fb:=\ker(\T\to\T'')$. Put $R:=\T\otimes_\Lam I_F$, $R':=\T'\otimes_\Lam I_F$,
$R'':=\T''\otimes_\Lam I_F$, and keep the names $\fa,\fb$ for $\fa R,\fb R$. $\lambda:R\to I_F$ is the $I_F$-linear extension of
$\lambda_F$. Since $F$ is $q$-new, $\lambda$ factors as $\lambda=\lambda'\circ(R\to R')$. For a finite $A$-algebra $S$ and an
$A$-algebra map $\mu:S\to A$ to a domain, $\eta(S,\mu):=\mu(\Ann_S\ker\mu)$ (Definition~\ref{def:s06-gor}). Put
\[\eta_F:=\eta(R,\lambda),\qquad \eta^B_F:=\eta(R',\lambda'),\qquad (h_F):=\eta_F\,(\eta^B_F)^{-1}\ \text{ when both are principal.}\]
\status{definition} (The quotient $h_F$ makes sense only under (Gor$_R$), (Gor$_{R'}$) and $\eta^B_F\ne0$; then $h_F\in I_F$ is well
defined up to $I_F^\times$, Thm.~\ref{thm:s06-E}.)

\medskip
\noindent\rptnum{20.5}\label{rpt:def:s06-ctam}\textbf{Definition 20.5 (combinatorial Tamagawa ideal).}
$\cc:=\mathrm{im}(\fb\to R')=(\fa+\fb)/\fa$, so $R'/\cc\cong R/(\fa+\fb)$ is the $\Lam$-adic old/new congruence module at $q$
(after $\otimes I_F$). Put $\cc_q(F):=\lambda'(\cc)=\lambda(\fb)\subseteq I_F$. At an arithmetic point $P$ with $q$-new
specialisation $f$, the \emph{fixed-weight} version is $\cc_q(f):=\lambda_f(\fb_k)$, where $\fb_k=\ker(\T_k(N)_\m\to\T''_k)$ is taken
in the weight-$k$ Hecke algebra. \status{definition}
(This is Diamond's $\Lam$-adic old/new congruence module \cite{Dia89}, read here along the new branch; \cite[Thm.~6A--6C, p.~80]{Dia89}
\stcited{} computes it along the old branch.)

\medskip
\noindent\rptnum{20.6}\label{rpt:def:s06-tam}\textbf{Definition 20.6 (Tamagawa ideal).}
Let $H$ be Hida's ordinary $\Lam$-adic étale cohomology of tame level $N$ localised at $\m$. It is free of rank two over $\T$
under (H2) (input (1) of Prop.~\ref{prop:s06-galois}). Put $\mathcal T_F:=H\otimes_{\T,\lambda_F}I_F$, with
$\rho_F:G_\Q\to\GL_2(I_F)$. Define $\tq_q(F)\subseteq I_F$ as the ideal generated by the entries of $\rho_F(\sigma)-1$,
$\sigma\in I_q$. The ideal does not depend on the basis of $\mathcal T_F$. \status{definition}
